
%%%%%%%%%%%%%%%%%%%%%%%%%%%%%%%%%%%%%%%%%
% The Legrand Orange Book
% LaTeX Template
% Version 2.0 (9/2/15)
%
% This template has been downloaded from:
% http://www.LaTeXTemplates.com
%
% Mathias Legrand (legrand.mathias@gmail.com) with modifications by:
% Vel (vel@latextemplates.com)
%
% License:
% CC BY-NC-SA 3.0 (http://creativecommons.org/licenses/by-nc-sa/3.0/)
%
% Compiling this template:
% This template uses biber for its bibliography and makeindex for its index.
% When you first open the template, compile it from the command line with the 
% commands below to make sure your LaTeX distribution is configured correctly:
%
% 1) pdflatex main
% 2) makeindex main.idx -s StyleInd.ist
% 3) biber main
% 4) pdflatex main x 2
%
% After this, when you wish to update the bibliography/index use the appropriate
% command above and make sure to compile with pdflatex several times 
% afterwards to propagate your changes to the document.
%
% This template also uses a number of packages which may need to be
% updated to the newest versions for the template to compile. It is strongly
% recommended you update your LaTeX distribution if you have any
% compilation errors.
%
% Important note:
% Chapter heading images should have a 2:1 width:height ratio,
% e.g. 920px width and 460px height.
%
%%%%%%%%%%%%%%%%%%%%%%%%%%%%%%%%%%%%%%%%%

%----------------------------------------------------------------------------------------
%	PACKAGES AND OTHER DOCUMENT CONFIGURATIONS
%----------------------------------------------------------------------------------------

\documentclass[11pt,fleqn]{book} % Default font size and left-justified equations


\usepackage{mathrsfs}
\usepackage{amsbsy}
\usepackage{graphicx}



\newcommand{\indep}{\rotatebox[origin=c]{90}{$\models$}}

%----------------------------------------------------------------------------------------

%fdlkjadlkjdfsal;kjafds;lkjafsd;lkjfsad

%%%%%%%%%%%%%%%%%%%%%%%%%%%%%%%%%%%%%%%%%
% The Legrand Orange Book
% Structural Definitions File
% Version 2.0 (9/2/15)
%
% Original author:
% Mathias Legrand (legrand.mathias@gmail.com) with modifications by:
% Vel (vel@latextemplates.com)
% 
% This file has been downloaded from:
% http://www.LaTeXTemplates.com
%
% License:
% CC BY-NC-SA 3.0 (http://creativecommons.org/licenses/by-nc-sa/3.0/)
%
%%%%%%%%%%%%%%%%%%%%%%%%%%%%%%%%%%%%%%%%%

%----------------------------------------------------------------------------------------
%	VARIOUS REQUIRED PACKAGES AND CONFIGURATIONS
%----------------------------------------------------------------------------------------

\usepackage[top=3cm,bottom=3cm,left=3cm,right=3cm,headsep=10pt,a4paper]{geometry} % Page margins

\usepackage{graphicx} % Required for including pictures
\graphicspath{{Pictures/}} % Specifies the directory where pictures are stored

\usepackage{lipsum} % Inserts dummy text

\usepackage{tikz} % Required for drawing custom shapes

\usepackage[english]{babel} % English language/hyphenation

\usepackage{enumitem} % Customize lists
\setlist{nolistsep} % Reduce spacing between bullet points and numbered lists

\usepackage{booktabs} % Required for nicer horizontal rules in tables

\usepackage{xcolor} % Required for specifying colors by name
\definecolor{ocre}{RGB}{243,102,25} % Define the orange color used for highlighting throughout the book

%----------------------------------------------------------------------------------------
%	FONTS
%----------------------------------------------------------------------------------------

\usepackage{avant} % Use the Avantgarde font for headings
%\usepackage{times} % Use the Times font for headings
\usepackage{mathptmx} % Use the Adobe Times Roman as the default text font together with math symbols from the Sym­bol, Chancery and Com­puter Modern fonts

\usepackage{microtype} % Slightly tweak font spacing for aesthetics
\usepackage[utf8]{inputenc} % Required for including letters with accents
\usepackage[T1]{fontenc} % Use 8-bit encoding that has 256 glyphs

%----------------------------------------------------------------------------------------
%	BIBLIOGRAPHY AND INDEX
%----------------------------------------------------------------------------------------

\usepackage[style=alphabetic,citestyle=numeric,sorting=nyt,sortcites=true,autopunct=true,babel=hyphen,hyperref=true,abbreviate=false,backref=true,backend=biber]{biblatex}
\addbibresource{bibliography.bib} % BibTeX bibliography file
\defbibheading{bibempty}{}

\usepackage{calc} % For simpler calculation - used for spacing the index letter headings correctly
\usepackage{makeidx} % Required to make an index
\makeindex % Tells LaTeX to create the files required for indexing

%----------------------------------------------------------------------------------------
%	MAIN TABLE OF CONTENTS
%----------------------------------------------------------------------------------------

\usepackage{titletoc} % Required for manipulating the table of contents

\contentsmargin{0cm} % Removes the default margin

% Part text styling
\titlecontents{part}[0cm]
{\addvspace{20pt}\centering\large\bfseries}
{}
{}
{}

% Chapter text styling
\titlecontents{chapter}[1.25cm] % Indentation
{\addvspace{12pt}\large\sffamily\bfseries} % Spacing and font options for chapters
{\color{ocre!60}\contentslabel[\Large\thecontentslabel]{1.25cm}\color{ocre}} % Chapter number
{\color{ocre}}  
{\color{ocre!60}\normalsize\;\titlerule*[.5pc]{.}\;\thecontentspage} % Page number

% Section text styling
\titlecontents{section}[1.25cm] % Indentation
{\addvspace{3pt}\sffamily\bfseries} % Spacing and font options for sections
{\contentslabel[\thecontentslabel]{1.25cm}} % Section number
{}
{\hfill\color{black}\thecontentspage} % Page number
[]

% Subsection text styling
\titlecontents{subsection}[1.25cm] % Indentation
{\addvspace{1pt}\sffamily\small} % Spacing and font options for subsections
{\contentslabel[\thecontentslabel]{1.25cm}} % Subsection number
{}
{\ \titlerule*[.5pc]{.}\;\thecontentspage} % Page number
[]

% List of figures
\titlecontents{figure}[0em]
{\addvspace{-5pt}\sffamily}
{\thecontentslabel\hspace*{1em}}
{}
{\ \titlerule*[.5pc]{.}\;\thecontentspage}
[]

% List of tables
\titlecontents{table}[0em]
{\addvspace{-5pt}\sffamily}
{\thecontentslabel\hspace*{1em}}
{}
{\ \titlerule*[.5pc]{.}\;\thecontentspage}
[]

%----------------------------------------------------------------------------------------
%	MINI TABLE OF CONTENTS IN PART HEADS
%----------------------------------------------------------------------------------------

% Chapter text styling
\titlecontents{lchapter}[0em] % Indenting
{\addvspace{15pt}\large\sffamily\bfseries} % Spacing and font options for chapters
{\color{ocre}\contentslabel[\Large\thecontentslabel]{1.25cm}\color{ocre}} % Chapter number
{}  
{\color{ocre}\normalsize\sffamily\bfseries\;\titlerule*[.5pc]{.}\;\thecontentspage} % Page number

% Section text styling
\titlecontents{lsection}[0em] % Indenting
{\sffamily\small} % Spacing and font options for sections
{\contentslabel[\thecontentslabel]{1.25cm}} % Section number
{}
{}

% Subsection text styling
\titlecontents{lsubsection}[.5em] % Indentation
{\normalfont\footnotesize\sffamily} % Font settings
{}
{}
{}

%----------------------------------------------------------------------------------------
%	PAGE HEADERS
%----------------------------------------------------------------------------------------

\usepackage{fancyhdr} % Required for header and footer configuration

\pagestyle{fancy}
\renewcommand{\chaptermark}[1]{\markboth{\sffamily\normalsize\bfseries\chaptername\ \thechapter.\ #1}{}} % Chapter text font settings
\renewcommand{\sectionmark}[1]{\markright{\sffamily\normalsize\thesection\hspace{5pt}#1}{}} % Section text font settings
\fancyhf{} \fancyhead[LE,RO]{\sffamily\normalsize\thepage} % Font setting for the page number in the header
\fancyhead[LO]{\rightmark} % Print the nearest section name on the left side of odd pages
\fancyhead[RE]{\leftmark} % Print the current chapter name on the right side of even pages
\renewcommand{\headrulewidth}{0.5pt} % Width of the rule under the header
\addtolength{\headheight}{2.5pt} % Increase the spacing around the header slightly
\renewcommand{\footrulewidth}{0pt} % Removes the rule in the footer
\fancypagestyle{plain}{\fancyhead{}\renewcommand{\headrulewidth}{0pt}} % Style for when a plain pagestyle is specified

% Removes the header from odd empty pages at the end of chapters
\makeatletter
\renewcommand{\cleardoublepage}{
\clearpage\ifodd\c@page\else
\hbox{}
\vspace*{\fill}
\thispagestyle{empty}
\newpage
\fi}

%----------------------------------------------------------------------------------------
%	THEOREM STYLES
%----------------------------------------------------------------------------------------

\usepackage{amsmath,amsfonts,amssymb,amsthm} % For math equations, theorems, symbols, etc

\newcommand{\intoo}[2]{\mathopen{]}#1\,;#2\mathclose{[}}
\newcommand{\ud}{\mathop{\mathrm{{}d}}\mathopen{}}
\newcommand{\intff}[2]{\mathopen{[}#1\,;#2\mathclose{]}}
\newtheorem{notation}{Notation}[chapter]

% Boxed/framed environments
\newtheoremstyle{ocrenumbox}% % Theorem style name
{0pt}% Space above
{0pt}% Space below
{\normalfont}% % Body font
{}% Indent amount
{\small\bf\sffamily\color{ocre}}% % Theorem head font
{\;}% Punctuation after theorem head
{0.25em}% Space after theorem head
{\small\sffamily\color{ocre}\thmname{#1}\nobreakspace\thmnumber{\@ifnotempty{#1}{}\@upn{#2}}% Theorem text (e.g. Theorem 2.1)
\thmnote{\nobreakspace\the\thm@notefont\sffamily\bfseries\color{black}---\nobreakspace#3.}} % Optional theorem note
\renewcommand{\qedsymbol}{$\blacksquare$}% Optional qed square

\newtheoremstyle{blacknumex}% Theorem style name
{5pt}% Space above
{5pt}% Space below
{\normalfont}% Body font
{} % Indent amount
{\small\bf\sffamily}% Theorem head font
{\;}% Punctuation after theorem head
{0.25em}% Space after theorem head
{\small\sffamily{\tiny\ensuremath{\blacksquare}}\nobreakspace\thmname{#1}\nobreakspace\thmnumber{\@ifnotempty{#1}{}\@upn{#2}}% Theorem text (e.g. Theorem 2.1)
\thmnote{\nobreakspace\the\thm@notefont\sffamily\bfseries---\nobreakspace#3.}}% Optional theorem note

\newtheoremstyle{blacknumbox} % Theorem style name
{0pt}% Space above
{0pt}% Space below
{\normalfont}% Body font
{}% Indent amount
{\small\bf\sffamily}% Theorem head font
{\;}% Punctuation after theorem head
{0.25em}% Space after theorem head
{\small\sffamily\thmname{#1}\nobreakspace\thmnumber{\@ifnotempty{#1}{}\@upn{#2}}% Theorem text (e.g. Theorem 2.1)
\thmnote{\nobreakspace\the\thm@notefont\sffamily\bfseries---\nobreakspace#3.}}% Optional theorem note

% Non-boxed/non-framed environments
\newtheoremstyle{ocrenum}% % Theorem style name
{5pt}% Space above
{5pt}% Space below
{\normalfont}% % Body font
{}% Indent amount
{\small\bf\sffamily\color{ocre}}% % Theorem head font
{\;}% Punctuation after theorem head
{0.25em}% Space after theorem head
{\small\sffamily\color{ocre}\thmname{#1}\nobreakspace\thmnumber{\@ifnotempty{#1}{}\@upn{#2}}% Theorem text (e.g. Theorem 2.1)
\thmnote{\nobreakspace\the\thm@notefont\sffamily\bfseries\color{black}---\nobreakspace#3.}} % Optional theorem note
\renewcommand{\qedsymbol}{$\blacksquare$}% Optional qed square
\makeatother

% Defines the theorem text style for each type of theorem to one of the three styles above
\newcounter{dummy} 
\numberwithin{dummy}{section}
\theoremstyle{ocrenumbox}
\newtheorem{theoremeT}[dummy]{Theorem}
\newtheorem{problem}{Problem}[chapter]
\newtheorem{exerciseT}{Exercise}[chapter]
\theoremstyle{blacknumex}
\newtheorem{exampleT}{Example}[chapter]
\theoremstyle{blacknumbox}
\newtheorem{vocabulary}{Vocabulary}[chapter]
\newtheorem{definitionT}{Definition}[section]
\newtheorem{corollaryT}[dummy]{Corollary}
\theoremstyle{ocrenum}
\newtheorem{proposition}[dummy]{Proposition}

%----------------------------------------------------------------------------------------
%	DEFINITION OF COLORED BOXES
%----------------------------------------------------------------------------------------

\RequirePackage[framemethod=default]{mdframed} % Required for creating the theorem, definition, exercise and corollary boxes

% Theorem box
\newmdenv[skipabove=7pt,
skipbelow=7pt,
backgroundcolor=black!5,
linecolor=ocre,
innerleftmargin=5pt,
innerrightmargin=5pt,
innertopmargin=5pt,
leftmargin=0cm,
rightmargin=0cm,
innerbottommargin=5pt]{tBox}

% Exercise box	  
\newmdenv[skipabove=7pt,
skipbelow=7pt,
rightline=false,
leftline=true,
topline=false,
bottomline=false,
backgroundcolor=ocre!10,
linecolor=ocre,
innerleftmargin=5pt,
innerrightmargin=5pt,
innertopmargin=5pt,
innerbottommargin=5pt,
leftmargin=0cm,
rightmargin=0cm,
linewidth=4pt]{eBox}	

% Definition box
\newmdenv[skipabove=7pt,
skipbelow=7pt,
rightline=false,
leftline=true,
topline=false,
bottomline=false,
linecolor=ocre,
innerleftmargin=5pt,
innerrightmargin=5pt,
innertopmargin=0pt,
leftmargin=0cm,
rightmargin=0cm,
linewidth=4pt,
innerbottommargin=0pt]{dBox}	

% Corollary box
\newmdenv[skipabove=7pt,
skipbelow=7pt,
rightline=false,
leftline=true,
topline=false,
bottomline=false,
linecolor=gray,
backgroundcolor=black!5,
innerleftmargin=5pt,
innerrightmargin=5pt,
innertopmargin=5pt,
leftmargin=0cm,
rightmargin=0cm,
linewidth=4pt,
innerbottommargin=5pt]{cBox}

% Creates an environment for each type of theorem and assigns it a theorem text style from the "Theorem Styles" section above and a colored box from above
\newenvironment{theorem}{\begin{tBox}\begin{theoremeT}}{\end{theoremeT}\end{tBox}}
\newenvironment{exercise}{\begin{eBox}\begin{exerciseT}}{\hfill{\color{ocre}\tiny\ensuremath{\blacksquare}}\end{exerciseT}\end{eBox}}				  
\newenvironment{definition}{\begin{dBox}\begin{definitionT}}{\end{definitionT}\end{dBox}}	
\newenvironment{example}{\begin{exampleT}}{\hfill{\tiny\ensuremath{\blacksquare}}\end{exampleT}}		
\newenvironment{corollary}{\begin{cBox}\begin{corollaryT}}{\end{corollaryT}\end{cBox}}	

%----------------------------------------------------------------------------------------
%	REMARK ENVIRONMENT
%----------------------------------------------------------------------------------------

\newenvironment{remark}{\par\vspace{10pt}\small % Vertical white space above the remark and smaller font size
\begin{list}{}{
\leftmargin=35pt % Indentation on the left
\rightmargin=25pt}\item\ignorespaces % Indentation on the right
\makebox[-2.5pt]{\begin{tikzpicture}[overlay]
\node[draw=ocre!60,line width=1pt,circle,fill=ocre!25,font=\sffamily\bfseries,inner sep=2pt,outer sep=0pt] at (-15pt,0pt){\textcolor{ocre}{R}};\end{tikzpicture}} % Orange R in a circle
\advance\baselineskip -1pt}{\end{list}\vskip5pt} % Tighter line spacing and white space after remark

%----------------------------------------------------------------------------------------
%	SECTION NUMBERING IN THE MARGIN
%----------------------------------------------------------------------------------------

\makeatletter
\renewcommand{\@seccntformat}[1]{\llap{\textcolor{ocre}{\csname the#1\endcsname}\hspace{1em}}}                    
\renewcommand{\section}{\@startsection{section}{1}{\z@}
{-4ex \@plus -1ex \@minus -.4ex}
{1ex \@plus.2ex }
{\normalfont\large\sffamily\bfseries}}
\renewcommand{\subsection}{\@startsection {subsection}{2}{\z@}
{-3ex \@plus -0.1ex \@minus -.4ex}
{0.5ex \@plus.2ex }
{\normalfont\sffamily\bfseries}}
\renewcommand{\subsubsection}{\@startsection {subsubsection}{3}{\z@}
{-2ex \@plus -0.1ex \@minus -.2ex}
{.2ex \@plus.2ex }
{\normalfont\small\sffamily\bfseries}}                        
\renewcommand\paragraph{\@startsection{paragraph}{4}{\z@}
{-2ex \@plus-.2ex \@minus .2ex}
{.1ex}
{\normalfont\small\sffamily\bfseries}}

%----------------------------------------------------------------------------------------
%	PART HEADINGS
%----------------------------------------------------------------------------------------

% numbered part in the table of contents
\newcommand{\@mypartnumtocformat}[2]{%
\setlength\fboxsep{0pt}%
\noindent\colorbox{ocre!20}{\strut\parbox[c][.7cm]{\ecart}{\color{ocre!70}\Large\sffamily\bfseries\centering#1}}\hskip\esp\colorbox{ocre!40}{\strut\parbox[c][.7cm]{\linewidth-\ecart-\esp}{\Large\sffamily\centering#2}}}%
%%%%%%%%%%%%%%%%%%%%%%%%%%%%%%%%%%
% unnumbered part in the table of contents
\newcommand{\@myparttocformat}[1]{%
\setlength\fboxsep{0pt}%
\noindent\colorbox{ocre!40}{\strut\parbox[c][.7cm]{\linewidth}{\Large\sffamily\centering#1}}}%
%%%%%%%%%%%%%%%%%%%%%%%%%%%%%%%%%%
\newlength\esp
\setlength\esp{4pt}
\newlength\ecart
\setlength\ecart{1.2cm-\esp}
\newcommand{\thepartimage}{}%
\newcommand{\partimage}[1]{\renewcommand{\thepartimage}{#1}}%
\def\@part[#1]#2{%
\ifnum \c@secnumdepth >-2\relax%
\refstepcounter{part}%
\addcontentsline{toc}{part}{\texorpdfstring{\protect\@mypartnumtocformat{\thepart}{#1}}{\partname~\thepart\ ---\ #1}}
\else%
\addcontentsline{toc}{part}{\texorpdfstring{\protect\@myparttocformat{#1}}{#1}}%
\fi%
\startcontents%
\markboth{}{}%
{\thispagestyle{empty}%
\begin{tikzpicture}[remember picture,overlay]%
\node at (current page.north west){\begin{tikzpicture}[remember picture,overlay]%	
\fill[ocre!20](0cm,0cm) rectangle (\paperwidth,-\paperheight);
\node[anchor=north] at (4cm,-3.25cm){\color{ocre!40}\fontsize{220}{100}\sffamily\bfseries\@Roman\c@part}; 
\node[anchor=south east] at (\paperwidth-1cm,-\paperheight+1cm){\parbox[t][][t]{8.5cm}{
\printcontents{l}{0}{\setcounter{tocdepth}{1}}%
}};
\node[anchor=north east] at (\paperwidth-1.5cm,-3.25cm){\parbox[t][][t]{15cm}{\strut\raggedleft\color{white}\fontsize{30}{30}\sffamily\bfseries#2}};
\end{tikzpicture}};
\end{tikzpicture}}%
\@endpart}
\def\@spart#1{%
\startcontents%
\phantomsection
{\thispagestyle{empty}%
\begin{tikzpicture}[remember picture,overlay]%
\node at (current page.north west){\begin{tikzpicture}[remember picture,overlay]%	
\fill[ocre!20](0cm,0cm) rectangle (\paperwidth,-\paperheight);
\node[anchor=north east] at (\paperwidth-1.5cm,-3.25cm){\parbox[t][][t]{15cm}{\strut\raggedleft\color{white}\fontsize{30}{30}\sffamily\bfseries#1}};
\end{tikzpicture}};
\end{tikzpicture}}
\addcontentsline{toc}{part}{\texorpdfstring{%
\setlength\fboxsep{0pt}%
\noindent\protect\colorbox{ocre!40}{\strut\protect\parbox[c][.7cm]{\linewidth}{\Large\sffamily\protect\centering #1\quad\mbox{}}}}{#1}}%
\@endpart}
\def\@endpart{\vfil\newpage
\if@twoside
\if@openright
\null
\thispagestyle{empty}%
\newpage
\fi
\fi
\if@tempswa
\twocolumn
\fi}

%----------------------------------------------------------------------------------------
%	CHAPTER HEADINGS
%----------------------------------------------------------------------------------------

\newcommand{\thechapterimage}{}%
\newcommand{\chapterimage}[1]{\renewcommand{\thechapterimage}{#1}}%
\def\@makechapterhead#1{%
{\parindent \z@ \raggedright \normalfont
\ifnum \c@secnumdepth >\m@ne
\if@mainmatter
\begin{tikzpicture}[remember picture,overlay]
\node at (current page.north west)
{\begin{tikzpicture}[remember picture,overlay]
\node[anchor=north west,inner sep=0pt] at (0,0) {\includegraphics[width=\paperwidth]{\thechapterimage}};
\draw[anchor=west] (\Gm@lmargin,-9cm) node [line width=2pt,rounded corners=15pt,draw=ocre,fill=white,fill opacity=0.5,inner sep=15pt]{\strut\makebox[22cm]{}};
\draw[anchor=west] (\Gm@lmargin+.3cm,-9cm) node {\huge\sffamily\bfseries\color{black}\thechapter. #1\strut};
\end{tikzpicture}};
\end{tikzpicture}
\else
\begin{tikzpicture}[remember picture,overlay]
\node at (current page.north west)
{\begin{tikzpicture}[remember picture,overlay]
\node[anchor=north west,inner sep=0pt] at (0,0) {\includegraphics[width=\paperwidth]{\thechapterimage}};
\draw[anchor=west] (\Gm@lmargin,-9cm) node [line width=2pt,rounded corners=15pt,draw=ocre,fill=white,fill opacity=0.5,inner sep=15pt]{\strut\makebox[22cm]{}};
\draw[anchor=west] (\Gm@lmargin+.3cm,-9cm) node {\huge\sffamily\bfseries\color{black}#1\strut};
\end{tikzpicture}};
\end{tikzpicture}
\fi\fi\par\vspace*{270\p@}}}

%-------------------------------------------

\def\@makeschapterhead#1{%
\begin{tikzpicture}[remember picture,overlay]
\node at (current page.north west)
{\begin{tikzpicture}[remember picture,overlay]
\node[anchor=north west,inner sep=0pt] at (0,0) {\includegraphics[width=\paperwidth]{\thechapterimage}};
\draw[anchor=west] (\Gm@lmargin,-9cm) node [line width=2pt,rounded corners=15pt,draw=ocre,fill=white,fill opacity=0.5,inner sep=15pt]{\strut\makebox[22cm]{}};
\draw[anchor=west] (\Gm@lmargin+.3cm,-9cm) node {\huge\sffamily\bfseries\color{black}#1\strut};
\end{tikzpicture}};
\end{tikzpicture}
\par\vspace*{270\p@}}
\makeatother

%----------------------------------------------------------------------------------------
%	HYPERLINKS IN THE DOCUMENTS
%----------------------------------------------------------------------------------------

\usepackage{hyperref}
\hypersetup{hidelinks,backref=true,pagebackref=true,hyperindex=true,colorlinks=false,breaklinks=true,urlcolor= ocre,bookmarks=true,bookmarksopen=false,pdftitle={Title},pdfauthor={Author}}
\usepackage{bookmark}
\bookmarksetup{
open,
numbered,
addtohook={%
\ifnum\bookmarkget{level}=0 % chapter
\bookmarksetup{bold}%
\fi
\ifnum\bookmarkget{level}=-1 % part
\bookmarksetup{color=ocre,bold}%
\fi
}
}
 % Insert the commands.tex file which contains the majority of the structure behind the template

\begin{document}

%----------------------------------------------------------------------------------------
%	TITLE PAGE
%----------------------------------------------------------------------------------------

\begingroup
\thispagestyle{empty}
\begin{tikzpicture}[remember picture,overlay]
\coordinate [below=12cm] (midpoint) at (current page.north);
\node at (current page.north west)
{\begin{tikzpicture}[remember picture,overlay]
\node[anchor=north west,inner sep=0pt] at (0,0) {\includegraphics[width=\paperwidth]{background}}; % Background image
\draw[anchor=north] (midpoint) node [fill=ocre!30!white,fill opacity=0.6,text opacity=1,inner sep=1cm]{\Huge\centering\bfseries\sffamily\parbox[c][][t]{\paperwidth}{\centering  Advanced Statistical Inference\\[15pt] % Book title
{\Large STAT 561 - Advanced Statistical Inference}\\[20pt] % Subtitle
{\huge Dr. Bing Li}}}; % Author name
\end{tikzpicture}};
\end{tikzpicture}
\vfill
\endgroup

%----------------------------------------------------------------------------------------
%	COPYRIGHT PAGE
%----------------------------------------------------------------------------------------

\newpage
~\vfill
\thispagestyle{empty}

\noindent Copyright \copyright\ 2013 John Smith\\ % Copyright notice

\noindent \textsc{Published by Publisher}\\ % Publisher

\noindent \textsc{book-website.com}\\ % URL

\noindent Licensed under the Creative Commons Attribution-NonCommercial 3.0 Unported License (the ``License''). You may not use this file except in compliance with the License. You may obtain a copy of the License at \url{http://creativecommons.org/licenses/by-nc/3.0}. Unless required by applicable law or agreed to in writing, software distributed under the License is distributed on an \textsc{``as is'' basis, without warranties or conditions of any kind}, either express or implied. See the License for the specific language governing permissions and limitations under the License.\\ % License information

\noindent \textit{First printing, March 2013} % Printing/edition date

%----------------------------------------------------------------------------------------
%	TABLE OF CONTENTS
%----------------------------------------------------------------------------------------

\chapterimage{chapter_head_1.pdf} % Table of contents heading image

\pagestyle{empty} % No headers

\tableofcontents % Print the table of contents itself

\cleardoublepage % Forces the first chapter to start on an odd page so it's on the right

\pagestyle{fancy} % Print headers again

%----------------------------------------------------------------------------------------
%	PART
%----------------------------------------------------------------------------------------

\part{Part One}

%----------------------------------------------------------------------------------------
%	CHAPTER 1
%----------------------------------------------------------------------------------------

\chapterimage{chapter_head_2.pdf} % Chapter heading image

\chapter{Basic Ideas in Bayesian Anaylysis}


% %------------------------------------------------

\textbf{Mathematical Preparation}\\

\textbf{Monday January 9}\\

\begin{enumerate}
	\item Product $\sigma$-Field\\

			$(\Omega_1, \mathcal{F}_1, \mu_1), (\Omega_2, \mathcal{F}_2, \mu_2)$ are two measure spaces. The goal is to construct a $\sigma$-field on $\Omega_1x\Omega_2$. \\

			Let $\mathcal{A} = \{AxB: A \in \mathcal{F}_1, B \in \mathcal{F}_2\}$. \\

			The $\sigma$-field generated $\mathcal{A}$ is called the product $\sigma$-field, written as $\mathcal{F}_1 x \mathcal{F}_2$, that is $\sigma(\mathcal{A})$. This is NOT a cartesian product, which would be $\{(A, B): A\in \mathcal{F}_1, B \in \mathcal{F}_2\}$. 

	\item Proctuct Measure\\

			Let $E \in \mathcal{F}_1 x \mathcal{F}_2$. Let $E_2(\omega_1) = \{\omega_2:(\omega_1, \omega_2) \in E\}$  and similarly, $E_1(\omega_2) = \{\omega_1:(\omega_1, \omega_2) \in E\}$. \\

			It is true (in Billingsly) that 

			\begin{theorem}[Number Unknown]
				If $E \in \mathcal{F}_1 x \mathcal{F}_2$ then $E_1(\omega_2) \in \mathcal{F}_1$ for all $\omega_2 \in \Omega_2$. Similarly, $E_2(\omega_1) \in \mathcal{F}_2$ for all $\omega_1 \in \Omega_1$. \\


				If $f: \Omega_1 x \Omega_2 \rightarrow \mathbb{R}$ measurable $\mathcal{F}_1 x \mathcal{F}_2 \setminus \mathcal{R}$. Then for each $\omega_1 \in \Omega_1$, 

						$$ f(\omega_1, \cdot) \textcircled{m} \mathcal{F}_2 \setminus \mathcal{R} \text{ for each } \omega_2 \in \Omega_2. $$

						$$ f(\cdot, \omega_2) \textcircled{m} \mathcal{F}_1 \setminus \mathcal{R} $$

				Now, for each $E \in \mathcal{F}_1 x \mathcal{F}_2$ consider 

						$$f_{1, E}: \Omega_1 \rightarrow \mathcal{R}, \omega_1 \mapsto \mu_2(E_2,(\omega_2)) $$

				It can be shown that $f_{1, E}$ is uniformly measurable $\mathcal{F}_1 \setminus \mathcal{R}$ for all E. 


			\end{theorem}

	\begin{proof}
		Outline. 

			\begin{itemize}
				\item Show that if $\mathcal{L} = \{E: f_{1, E} \textcircled{m} \mathcal{F}_1 \setminus \mathcal{R}\}$ then $\mathcal{L}$ is a $\lambda$-system. 
				\item Let $\mathcal{P} = \{AxB: A \in \mathcal{F}_1, B \in \mathcal{F}_2\}$ then it is a $\pi$-system.
			

				Furthermore, if $ E = A x B$, 
				
						$$E_2(\omega_1) = \left\{ \begin{array}{ll}
							B & \omega_1 \in A\\
							\emptyset & \omega_1 \notin A
						\end{array} \right. $$	

				So, $\mu_2(E_2(\omega_1)) = \left\{ \begin{array}{ll}
							\mu_2(B) & \omega_1 \in A\\
							\emptyset & \omega_1 \notin A
						\end{array} \right. = I_A(\omega_1)\mu(B) = f_{1, E}$

				So, $f_{1, E} \textcircled{m} \mathcal{F}_1$. 

				Thus $\mathcal{P} \subseteq \mathcal{L}$. 	
				
				\item By $\pi-\lambda$ Theorem, $\mathcal{F}_1 x \mathcal{F}_2 \subseteq \mathcal{L}$. 

	\end{itemize}

		Similarly, $f_{2, E} \textcircled{m} \mathcal{F}_2 \setminus \mathcal{R}$. 

		We can now define two set functions, 

				$$\pi'(E) = \int f_{1, E} d\mu_1 $$
				$$\pi''(E) = \int f_{2, E} d\mu_2 $$

		Again using $\pi-\lambda$ Theorem, it can be shown that, $\pi', \pi''$ are both measure and if $\mu_1, \mu_2$ are $\sigma$-finite, then

				$$\pi' = \pi'' \text{ on } \mathcal{F}_1 x \mathcal{F}_2 $$

		Note that here, $\mathcal{P}$ equals $\mathcal{A}$ used at begining	of notes.


		We did not have a measure in $\mathcal{F}_1 x \mathcal{F}_2$. Now we have $\pi', \pi''$ both measures on $\mathcal{F}_1 x \mathcal{F}_2$, they are the same. We call this measure the product meaure, written as $\mu_1 x \mu_2$. 


		Note that $(\Omega_1 x \Omega_2, \mathcal{F}_1 x \mathcal{F}_2, \mu_1 x \mu_2)$ is called product measure space. 
	\end{proof}

	\item Tonelli's Theorem\\

			$(\Omega_1, \mathcal{F}_1, \mu_1), (\Omega_2, \mathcal{F}_2, \mu_2)$ are two $\sigma$-finite measure spaces. \\

			$(\Omega_1 x \Omega_2, \mathcal{F}_1 x \mathcal{F}_2, \mu_1 x \mu_2)$ is the product measure space.\\

			Suppose we have $f:\Omega_1 x \Omega_2 \rightarrow \mathbb{R} \textcircled{m} \mathcal{F}_1 x \mathcal{F}_2 \setminus \mathcal{R} $. Where $f \geq 0$ and 

					$$\int f d (\mu_1 x \mu_2) =\int \left[ \int(f(\cdot, \omega_2) d\mu_1) \right] d\mu_2$$

	\item Fubini's Theorem\\

	The conclusion of Tonelli's Theorem still holds if f is NOT nonnegative, but if f is integrable $\mu_2$. (integrable - integral of absolute value of function is finite) \\

\textbf{Wednesday January 11}\\	

	\item Conditional Probability\\

	This is a special application of Radon- Nikodgm Theorem. We know that 

			$$P(A|B) = \frac{P(A, B)}{P(B)} $$

	We may define $P(A|\mathcal{G})$ when $\mathcal{G}\subseteq \mathcal{F}$ as sub-$\sigma$-field. We defined this intuitively in elementary probability course (definition above), but we ar enot going to define it generally. \\

	Now let $A \in \mathcal{F}$ and $\mathcal{G} \subset \mathcal{F}$ be a $\sigma$-field. Consider the set function 

			$$\nu: \mathcal{G} \rightarrow \mathbb{R}, G \mapsto P(AG)$$ 

	It can be easily shown that $\nu$ is a measure on $\mathcal{G}$. Consider another set function, 

			$$\mu: \mathcal{G} \rightarrow \mathbb{R}, G \mapsto P(G) $$

	So $\mu$ is nothting but P restricted on $\mathcal{G}$.\\

	It's easy to show that $\nu << \mu$. 

			$$\mu(G) = 0 \Rightarrow P(G) = 0 \Rightarrow P(AG) = 0 \Rightarrow \nu(G) = 0$$

		By Radon-Nikodgm Theorem, there exists a $\delta$ such that

			$$\nu(G) = \int_G \delta d\mu \quad \forall G \in \mathcal{G} $$

		$\delta$ is called R-N Derivative, written as 

				$$\delta = \frac{d\nu}{d\mu} $$

		and is similar in  to $\frac{P(AG)}{P(G)}$, but it's more general. \\

		$\delta$ is called the conditional probability of A given $\mathcal{G}$. To distinguish it form P(A|B), where B is a set, we use $P(A||\mathcal{G})$, where $\mathcal{G}$ is a $\sigma$-field. By construction, \\

				\begin{enumerate}
					\item $\delta$ is measurable $\mathcal{G}$
					\item $\int_G \delta dp = P(AG) \quad \forall G\in\mathcal{G}$ \\
				\end{enumerate}


		Note that, by RNT, $\delta$ is unique with probability 1. Any $\delta'$ satisfying (a) and (b) has $\delta' = \delta a.e. P$. So, we say that $\delta$ is a version of conditional probability. \\

		So, $\delta$ is a version of $P(A||\mathcal{G})$ if and only if (a) and (b) are satisfied. We may define $P(A||\mathcal{G})$ either by RNT or (a) and (b). \\


		Properties of Conditional Probability\\

		It behaves like probability, but since it is a function, unique up to a.e. P, these properties have to be qualified by a.s. P.

		\begin{enumerate}	
			\item $P(\emptyset || \mathcal{G}) = 0, P(\Omega||\mathcal{G}) = 1$ a.s. P
			\item $0 \leq P(A|| \mathcal{G}) \leq 1$ a.s. P
			\item If $A_1, A_2, \dots$ are disjoint members of $\mathcal{F}$ then $P(\bigcup_n A_n || \mathcal{G}) = \sum_n P(A_n || \mathcal{G})$ a.s. P
		\end{enumerate}

		Let's consider the special case where $\mathcal{G}$ is a $\sigma$-field generated by some random element, T (i.e. $\mathcal{G} = \sigma(T)$). More specifically, for some measurable space $(\Omega_T, \mathcal{F}_T)$ where 


				$$T: \Omega \rightarrow \Omega_T \textcircled{m} \mathcal{F}\setminus\mathcal{F}_T \quad \mathcal{G} = T^{-1}(\mathcal{F}_T)$$

		Here, we write

		\begin{align*}
			P(A || \mathcal{G}) &= P(A || \sigma(T))\\
				&=P(A|| T^{-1}(\mathcal{F}_T))\\
				&= P(A || T)
		\end{align*}

		The following theorem makes checking that something is a conditional probability easier. In principle, we have to check $\int_G \delta dp = P(AG) \quad \forall G\in\mathcal{G}$. 

		\begin{theorem}[33.1 in Billingsly]
			Let $\mathcal{P}$ be a $pi$-system generating $\mathcal{G}$ and suppose that $\Omega$ is a countable union of sets in $\mathcal{P}$. An integrable function, $f$, is a version of $P(A || \mathcal{G})$ if 

				\begin{enumerate}
					\item f is measurable $\mathcal{G}$
					\item $\int_G f dp = P(AG) \quad \forall G \in \mathcal{P}$ 
				\end{enumerate}
		\end{theorem}

	\item Conditional Distribution\\

	Let there be probability space $(\Omega, \mathcal{F}, P)$, measurable space $(\Omega_X, \mathcal{F}_X)$, and a random element, $X:\Omega \rightarrow \Omega_X \textcircled{m} \mathcal{F}\setminus\mathcal{F}_X$. Also, let $\mathcal{G} \subseteq \mathcal{F}$ be a sub $\sigma$-field. \\

	We are going to define conditional distribution of X given G. Under very mild conditions there is a function

			$$f: \mathcal{F}_X x \Omega \rightarrow \mathbb{R} $$

	such that for each $A \in \mathcal{F}_X$, $f(A, \cdot)$ is a version of 

			$$P(X \in A || \mathcal{G}) = P(X^{-1}(A) || \mathcal{G}) $$

	and, for each $\omega \in \Omega$,  $f(\cdot, \omega)$ is a probability measure on $(\Omega_X, \mathcal{F}_X)$. \\

	The only condition for this existance is $(\Omega_X, \mathcal{F}_X)$ must be a Borel Space, that is $\mathcal{F}_X$ is Borel $\sigma$-field. This should always be the case for our purposes. 
	

	\item Conditional Expectation\\

	Let us have the same probability space, measurable space, random element, and sub $\sigma$-field as defined before, but here with $\bar{\mathbb{R}}$.  \\

	We want to define conditional expectaiton of X given $\mathcal{G}$. \\

	First, assume $X \geq 0$. Consider a set function, 

			$$\nu: \mathcal{G} \rightarrow \mathbb{R}, G \mapsto \int_G X dP $$ 

	It can be easily shown that $\nu$ is a measure. 

	Let $\mu$ again be $\mathcal{G} \rightarrow \mathbb{R}, G \mapsto  P(G)$. Then $\nu << \mu$. By RNT, $\delta = \frac{d\nu}{d\mu}$ is well defined. This is defined to be conditional expectation of X given $\mathcal{G}$, written as 

			$$E(X|| \mathcal{G}) $$

	Suppose $X \ngeq 0$, but integrable P. Recall that $X = X^+ - X^-$. Since $X^+, X^- \geq 0$, then both $E(X^+|| \mathcal{G}), E(X^-|| \mathcal{G})$ are defined by RNT. We define, 

			$$E(X|| \mathcal{G}) = E(X^+|| \mathcal{G}) - E(X^-|| \mathcal{G}) $$

				 

\textbf{Friday January 13}

	As in the case of $P(A || \mathcal{G})$, the equivalent conditions for $d: \Omega \rightarrow \mathbb{R} $ is a version of $E(X||\mathcal{G})$. 

			\begin{enumerate}
				\item $\delta$ measurable $\mathcal{G}$
				\item $\int_G \delta dP = \int X dP \quad \forall G \in \mathcal{G}$ 
			\end{enumerate}

			INSERT PHOTO FROM BOARD - "Mesh"


	The value of $\delta$ in each thick outlined cell is the average (with respect to P measure ) of $X(\omega)$ over the subcells (thin outlined) in thick cells. 

	We see from this definition that if  $ A \in \mathcal{F}$, $X = I_A$ then the second condition becomes 

			$$\int_G \delta dP = \int_G I_A dP = P(A \bigcap G) $$

	So, $E(I_A||\mathcal{G}) = P(A || \mathcal{G})$. 

	\textbf{Properties of Conditional Expectations}\\

	\begin{theorem}{34.2 in Billingsly}
		Suppose that $X, Y, X_n$ are integrable P. 

			\begin{enumerate}
				\item If X = a a.e. P, then $E(X || \mathcal{G})$ a.s. P
				\item $a, b \in \mathbb{R}$ then

						$$E(aX + bY || \mathcal{G}) = a(E(X ||\mathcal{G})) + b(E(X || \mathcal{G})) a.s. P $$
				\item If $X \leq Y$ a.s. P then 

						$$E(X ||\mathcal{G}) \leq E(Y ||\mathcal{G}) $$
				\item $|E(X||\mathcal{G})| \leq E(|X| ||\mathcal{G})$ a.s. P (in fact this is true for all convex functions).
				\item If $X_n \rightarrow X$ a.s. P, $|X_n| \leq Y,$ and Y integrable P, then

						$$E(X_n ||\mathcal{G}) \rightarrow E(X ||\mathcal{G}) a.s. P $$
			\end{enumerate}

	\end{theorem}

	\begin{proof}
		Found in Billingsly. 
	\end{proof}

	\begin{theorem}[34.4 in Billingsly]
		If $\mathcal{G}_1 \subseteq \mathcal{G}_2 \subset \mathcal{F}$ and X integrable P, then

				$$E(E(X ||\mathcal{G}_2)|| \mathcal{G}_1) = E(X ||\mathcal{G}_1) $$

		This is called the Law of Iterative Conditional Expectation. 
	\end{theorem}

	\begin{theorem}[34.3 in Billingsly]
		
		If $X$ measurable $\mathcal{G}$, $Y \textcircled{m} \mathcal{F}$, then

				$$E(X Y||\mathcal{G}) = XE(Y ||\mathcal{G}) a.s. P $$
	\end{theorem}

	Other Properties

		\begin{enumerate}
			\item X, Y are random elements such that XY integrable P.
			\item If $\mathcal{G} \subseteq \mathcal{F}$ is the sub $\sigma$-field, then

					$$E(X E(Y||\mathcal{G})) = E(E(X ||\mathcal{G})Y) = E(E(X ||\mathcal{G}) E(Y ||\mathcal{G})) $$ 

			Conditional expectation is a self-adjoint operation. 

			\begin{proof}

			"Wire Theorem"\\

				$$\begin{aligned}
								E(X E(Y ||\mathcal{G}))	& = E( E(X E(Y ||\mathcal{G}) ||\mathcal{G}))\\
									&= E( E(Y ||\mathcal{G})E(X ||\mathcal{G}))	\\
									&= E( E(E(X ||\mathcal{G}) Y ||\mathcal{G})) \\
									&= 	E(E(X ||\mathcal{G}) Y)	
								\end{aligned}$$
			\end{proof}
		\end{enumerate}

	\item Conditional Distribution of a Random Element Given Another Random Element\\

		Here we have the typical probability space, measurable spaces for X and Y. 

		Let there be a function, 

			$h: \mathcal{F}_X x \Omega_Y \rightarrow \mathbb{R}$

		This function is called the conditional distribution of X given Y if 

				$$\tilde{h}(A, \omega) = h(A, Y(\omega)) $$

		We say that $\tilde{h}: \mathcal{F}_X x \Omega \rightarrow \mathbb{R}$ is the condiitional distribution of X given $\mathcal{G} = Y^{-1}(\mathcal{F}_Y$. 

		That is, 

				\begin{enumerate}
				 	\item For each $A \in \mathcal{F}_X$

				 			$$\tilde{h}(A, Y(\cdot)) = P(X^{-1}(A) || Y^{-1}(\mathcal{F}_Y)) $$
		 			\item  For each $\omega \in \Omega$

		 					$$\tilde{h}(\cdot, Y(\omega)) = P_{X|Y}(A | y) $$
				 \end{enumerate} 


	\item Conditional Density of One Random Element Given Another Random Element\\

		Suppose probability space and $\sigma$-finite measure spaces for X and Y. 

		Here our relevant function is

				$$g: \Omega_X x \Omega_Y $$

		which is the conditional density of X given Y if for all $A \in \mathcal{F}_X$, 

				$$\int_A g(x, y) d \mu_X(x) = P_{X|Y}(A|y) $$

		In the following special case, g ahs an explicit formula. 


		$(\Omega, \mathcal{F}, P)$\\
		$(\Omega_X, \mathcal{F}_X, \mu_X)$\\
		$(\Omega_Y, \mathcal{F}_Y, \mu_Y)$\\
		$(\Omega_X x \Omega_Y, \mathcal{F}_X x \mathcal{F}_Y, \mu_X x \mu_Y)$\\
		$(X, Y): \Omega \rightarrow \Omega_X x \Omega_Y \textcircled{m} \mathcal{F}\setminus \mathcal{F}_X x \mathcal{F}_Y$\\

		Let $P_X = P X^{-1}, P_Y = P Y^{-1}, P_{XY} = P (XY)^{-1}$. \\
		Assume $P_X << \mu_X, P_Y << \mu_Y, P_{XY} << \mu_X x \mu_Y$. 



				$$f_X = \frac{dP_X}{d\mu_X} $$
				$$f_Y = \frac{dP_Y}{d\mu_Y} $$ 
				$$f_{XY} = \frac{dP_{XY}}{d(\mu_X x \mu_Y) } $$ 

		Let

				$$f_{X|Y} = \begin{array}{ll}
				\frac{f_{XY}}{f_Y} & if f_Y \neq 0\\
				0 & f_Y = 0
					
				\end{array} $$

				$$f_{Y|X} = \begin{array}{ll}
				\frac{f_{XY}}{f_X} & if f_X \neq 0\\
				0 & f_X = 0
					
				\end{array} $$

		Then it is easy to show that each is indeed the conditional density of their respective elements (first given second). 
 
\textbf{Wednesday January 18}\\

Claim: $g(x, y)$ is the conditional density. 

\begin{proof}
	Want to show that for all $A \in \mathcal{F}_X$, 

			$$ \int_A g(x, y) d\mu_x(x) = P_{X|Y} (A|y) $$

	Which means that

			$$\int_A g(x, y(\omega)) d\mu_x(x) = P_{X|Y} (X^{-1}(A)|\sigma(y)) $$

	This is true if for all $G' \in \sigma(y)$

			$$\int_{G'} \int_A g(x, y(\omega)) d\mu_X (x) dP(\omega) = P(X^{-1}(A) \bigcap G') $$


	But note that 

				\begin{align*}
					G' &\in \sigma(y)\\
					\Leftrightarrow G' &\in Y^{-1}(\mathcal{F}_Y)\\
					G' &= Y^{-1}(G)\text{ for some } G\in \mathcal{F}_Y
				\end{align*}



	So we want to check that  

			$$\int_{Y^{-1}(G)} \int_A g(x, y(\omega)) d\mu_X (x) dP(\omega) = P(X^{-1}(A) \bigcap Y^{-1}(G)) $$

			\begin{align*}
				\int_{Y^{-1}(G)} \int_A g(x, y(\omega)) d\mu_X (x) dP(\omega) &= \int_{G} \int_A g(x, y) d\mu_X (x) dP_Y(y) \\
					&=\int_{G} \int_A \frac{f_{XY}(x, y)}{f_Y(y)} d\mu_X (x) [f_Y(y)] d\mu_Y(y) \\
					&=\int_{G} \int_A f_{XY}(x, y) d\mu_X (x) d\mu_Y(y)\\
					&=\int_{GxA} f_{XY}(x, y) d(\mu_X x \mu_Y)(x,y)\\
					&= P_{XY}(GXA) \\
					&= P \circ (X, Y)^{-1}(Ax G)\\
					&= P(X \in A, Y \in G)\\
					&= P(\omega: \omega \in X^{-1}(A) \& \omega \in Y^{-1}(G))\\
					&= P(X^{-1}(A) \bigcap Y^{-1}(G)) 
			\end{align*}
\end{proof}

\end{enumerate}
	
\section{Frequentist \& Bayesian Settings}\index{Frequentist \& Bayesian Settings}

We have our probability space $(\Omega, \mathcal{F}, P)$. We also have some data, 

		$$(\Omega_X, \mathcal{F}_X, \mu_X)$$ 
		$$ X: \Omega \rightarrow \Omega_X \textcircled{m} \mathcal{F}/ \mathcal{F}_X$$

Here, usually $\Omega_X$ is a $\mathbb{R}^m$.\\

Typically we have 

	$$X = (X_1, \dots, X_n)$$

and possibly, 

	$$ X_i = \begin{pmatrix}
		X_{i1}\\
		\vdots\\
		X_{ip}
	\end{pmatrix}$$

We could say that these data are independent and identically distributed (iid) random vectors of dimension p. In this case m = np. \\

The goal of statical inference is to estimate.

		$$P_X = PX^{-1} = P_0$$

The ?? distribution of X. \\

There are two schools of thought

	\begin{enumerate}
		\item Frequentist Approach - assume a family of distributions, $\mathcal{P}$, where $\mathcal{P} << \mu_x$. Usually we assume that $\mathcal{P}$ is a parametric family, $\mathcal{P} = \{P_\theta: \theta \in \Omega_\theta \subseteq \mathbb{R}^p \}$. We assume that $P_0 \in \mathcal{P}$, that is there exisgts $\theta_0 \in \Omega_\theta$ such that $P_\theta = P_0$. The goal is to estimate $P_0$. 
		\item Bayesian Approach - here we assume the data is generated by the conditional distribtuion $P_{X|\theta}$. We observe X, then determine what is the best estimate of the random $\theta.$
\end{enumerate}


\section{Prior Posterior \& Likelihood}


Here let there be probability space $(\Omega, \mathcal{F}, P)$; $\sigma$-finite measurable spaces $(\Omega_X, \mathcal{F}_X, \mu_X)$, $(\Omega_\theta, \mathcal{F}_\theta, \mu_\theta)$. Together, 

		$$(\Omega_X x \Omega_\theta, \mathcal{F}_X x \mathcal{F}_\theta, \mu_X x \mu_\theta)$$ 


Also, a random element, 

		$$(X, \theta):\Omega \rightarrow \Omega_X x \Omega_\theta \textcircled{m} \mathcal{F}\setminus\mathcal{F}_X x \mathcal{F}_\theta$$


$P_X = P \circ X^{-1} \leftarrow$ marginal distribution of X\\
$P_\theta = P \circ \theta^{-1} \leftarrow$ prior distribution \\
$P_{X, \theta} = P \circ (X, \theta)^{-1} \leftarrow$ joint distribution of X and $\theta$\\
$P_{X|\theta}(A|\theta): \mathcal{F}_X x \Omega_\theta \rightarrow \mathbb{R}$. Likeliehood distribution\\
$P_{\theta|X}(G|x): \mathcal{F}_\theta x \Omega_X \rightarrow \mathbb{R}$. Posterior distribution\\


Note in the following the first inequalities are \textbf{assumed}.\\

$P_X << \mu_X \Rightarrow f_X = \frac{d P_X}{d\mu_X}$ Marginal Density\\
$P_\theta << \mu_\theta \Rightarrow \pi_\theta = \frac{d P_\theta}{d\mu_\theta}$ Prior Density\\
$P_(X, \theta) << \mu_X x \mu_\theta \Rightarrow f_{X, \theta} (x, \theta) = \frac{d P_{X, \theta}}{d(\mu_x x \mu_\theta)}$ Joint Density\\


FINISH FROM PHOTO


One way to estimate $\theta$ is by maximizing $\pi_{\theta|X}(\theta|x)$. Want to do so with value that is most likely to happen (given the data). 

		$$\pi_{\theta|X} (\theta|x) = P_{\theta|X}(\theta = \theta | x) $$

By construction, 

		\begin{align*}
			\pi_{\theta|X} &= \frac{f_X\theta}{f_X}\\
				&= \frac{f_{X|\theta}\pi_\theta}{\int_{\Omega_\theta} f_{X|\theta} \pi_\theta d\mu_theta}
		\end{align*}


\section{Conditional Independence and Frquentist/Bayesian Sufficiency}\index{Conditional Independence and Frquentist/Bayesian Sufficiency}


\textbf{Independence}\\

Two random elements are said to be independent if for all $A' \in \sigma(X), G' \in \sigma(\theta)$ we have

		$$P(A' \bigcap G') = P(A')P(G') $$

This statement can also be expressed in $(\Omega_X x \Omega_\theta, \mathcal{F}_X x \mathcal{F}_\theta, P_X x P_\theta)$ as follows. \\

Since $A' \in \sigma(X) = X^{-1} (\mathcal{F}_X), A' = X^{-1}(A)$ for some $A \in \mathcal{F}_X$. So $G' = \Theta^{-1}(G). G \in \mathcal{F}_\Theta$. 

		\begin{align*}
			P(A' \bigcap G') &= P(X^{-1}(A) \bigcap \Theta^{-1}(G) )\\
				&=P(\{ \omega: \omega \in  X^{-1}(A) \bigcap \Theta^{-1}(G)\} )\\
				&=P(\{ \omega: \omega \in  X^{-1}(A) \& \omega \in \Theta^{-1}(G)\} )\\
				&=P(\{ \omega: X(\omega) \in A, \Theta(\omega) \in G\} )\\
				&=P(\{ \omega:( X(\omega),  \Theta(\omega)) \in A x G\} )\\
				&=P(\{ \omega:( X, \Theta)(\omega) \in A x G\} )\\
				&=P(\{ \omega:\omega \in( X, \Theta)^{-1} A x G\} )\\
				&=[P \circ ( X, \Theta)^{-1} ](A x G)\\
				&=P_{X, \Theta}(A x G)
		\end{align*}

Also note that

		$$P(A') = P(X^{-1}(A)) = P_X(A) $$
		$$P(G') = P_\Theta (G) $$


So with independence, (and for $A \in \mathcal{F}_X, G \in \mathcal{F}_\Theta$)

		$$P_{X, \Theta}(A x G) = P_X (A) P_\Theta(G) $$

But we know that this implies that $P_{X, \Theta}$ is the product measure $P_X x P_\Theta$. 

\textbf{Conditional Independence}\\

Now, given sub $\sigma$-field $\mathcal{G} \in \mathcal{F}$ we want to define $X \& \Theta$ conditionally independent given $\mathcal{G}$. 

\begin{definition}
	We say that $X \& \Theta$ are conditionally independent given $\mathcal{G}$ (i.e. $X \indep \Theta | \mathcal{G}$) if for all $A' \in \sigma(X), G' \in \sigma(\Theta)$ we have 

			$$P[A' \cap G' || \mathcal{G}] = P[A' || \mathcal{G}]P[G' || \mathcal{G}] a.s. P $$

	Equivalently for all $A \in \mathcal{F}_X, G \in \mathcal{F}_\Theta$, 

			$$P[X^{-1}(A) \cap \Theta^{-1}(G) || \mathcal{G}] =P[X^{-1}(A) || \mathcal{G}]P[\Theta^{-1}(G) || \mathcal{G}]  $$

	Equivalently, 

			$$P_{X, \Theta | \mathcal{G}} (A x G | \mathcal{G}) = P_{X|\mathcal{G}}(A | \mathcal{G}) P_{\Theta|\mathcal{G}} (G | \mathcal{G}) $$
\end{definition}


\textbf{Equivalent Condition for Conditional Independence}\\

\begin{theorem}[1.1 in Notes]
	The following statements are equivalent. 

	\begin{enumerate}
		\item $X \indep \Theta | \mathcal{G}$
		\item $P(X^{-1}(A) || \Theta, \mathcal{G}) = P(X^{-1}(A)| \mathcal{G}) a.s. P \quad \forall A \in \sigma(X)$
		\item $P(\Theta^{-1}(G) || X, \mathcal{G}) = P(\Theta^{-1}(G)|| \mathcal{G}) a.s. P \quad \forall G \in \sigma(\Theta)$
	\end{enumerate}
\end{theorem}

\begin{proof}
	It suffies to proof that $1 \Leftrightarrow 2$.\\

	$1 \Rightarrow 2$. We know that for all $A \in \mathcal{F}_X, G \in \mathcal{F}_\Theta$ that

			$$ P[X^{-1}(A) \cap \Theta^{-1}(G) || \mathcal{G}] =P[X^{-1}(A) || \mathcal{G}]P[\Theta^{-1}(G) || \mathcal{G}]$$

	Want that for all $A \in \mathcal{F}_X$ that $P(X^{-1}(A) || \Theta, \mathcal{G}) = P(X^{-1}(A)| \mathcal{G})$. 

	\begin{align*}
		P(X^{-1}(A) || \Theta, \mathcal{G}) &\equiv P\left(X^{-1}(A)|| \sigma(\sigma(\Theta) \cup \mathcal{G})\right) \\
			&= P(\dots || \sigma(\Theta^{-1}(\mathcal{F}_\Theta) \cup \mathcal{G}))
	\end{align*}

	So it suffices to show that 

			$$P(X^{-1}(A) || \sigma(\Theta^{-1}(\mathcal{F}_\Theta) \cup \mathcal{G})) = P(X^{-1}(A) || \mathcal{G})$$

	From the definition given we want to show that the above statement is true. which is so that the for all $B \in \sigma(\Theta^{-1}(\mathcal{F}_\Theta) \cup \mathcal{G} )$, 

			$$\int_B P(X^{-1}(A) || \mathcal{G}) dP = P(X^{-1} (A) \cap B) $$

	But this is very hard because B is hard to characterize. But we have theorem that says you only have to check (*) for all B in a $\pi$-system generating $\sigma(\Theta^{-1}(\mathcal{F}_\Theta)\cup \mathcal{G})$.

			$$\mathcal{P} = \{\Theta^{-1} (G) \cap F: G \in \mathcal{F}_\Theta, F \in \mathcal{G} \} $$

	It is trivial to show that $\mathcal{P}$ is a $\pi$-system. 

	MORE IN PHOTO

	Meanwhile, 

			$$\mathcal{P} \subseteq \sigma(\Theta^{-1} (\mathcal{F}_\Theta) \cup \mathcal{G}) $$

	Therefore, 

			$$\sigma(\Theta^{-1}(\mathcal{F}_\Theta)\cup \mathcal{G}) = \sigma(\mathcal{P}) $$

	So, sufficent to check (*) $\forall B \in \mathcal{P}'$

			$$B \in \mathcal{P} \Rightarrow B = \Theta^{-1}(G) \cap F, G \in \mathcal{F}_\Theta, F \in \mathcal{G} $$

	So, we want

			$$\int_{\Theta^{-1}(G)\cap F} P(X^{-1}(A) || \mathcal{G}) dP = P(\Theta^{-1}(G) \cap F \cap X^{-1}(A)) $$


			\begin{align*}
				\int_{\Theta^{-1}(G)\cap F} P(X^{-1}(A) || \mathcal{G}) dP &= \int_{\Theta^{-1}(G) \cap F} E \left(I_{X^{-1}(A)} || \mathcal{G} \right) dP\\
						&= E\left(I_{\Theta^{-1}(G)} I_F E(I_{X^{-1}(A)} || \mathcal{G})\right)\\
						&= E\left( E(I_{\Theta^{-1}(G)} I_F || \mathcal{G}) E(I_{X^{-1}(A)} || \mathcal{G})\right)\\
						&= E\left( I_F E(I_{\Theta^{-1}(G)}  || \mathcal{G}) E(I_{X^{-1}(A)} || \mathcal{G})\right)\\
						&= E\left( I_F E(I_{\Theta^{-1}(G)} I_{X^{-1}(A)} || \mathcal{G})\right)\\
						&= E\left(  E(I_F I_{\Theta^{-1}(G)} I_{X^{-1}(A)} || \mathcal{G})\right)\\
						&= E\left( I_F I_{\Theta^{-1}(G)} I_{X^{-1}(A)}\right)\\
						&= P(F \cap \Theta^{-1}(G) \cap X^{-1}(A))
			\end{align*}

\textbf{Monday January 23}\\

	$2 \Rightarrow 1$. We want to show that 

			$$P(X^{-1}(A)|| \mathcal{G}) P(\Theta^{-1}(G) || \mathcal{G}) $$

		is conditional probability of 

			$$P(X^{-1}(A) \cap \Theta^{-1}(G)|| \mathcal{G}) $$

		for all $F \in \mathcal{G}$. 
		
				\begin{align*}
						\int_F P(X^{-1}(A)|| \mathcal{G})P(\Theta^{-1}(G) || \mathcal{G}) dP &= E\left[I_F E\left(I_{X^{-1}(A)} || \mathcal{G} \right) E\left(I_{\Theta^{-1}(G) }|| \mathcal{G} \right) \right]\\
								&= E\left[ E\left(I_{X^{-1}(A)} || \mathcal{G} \right) E\left(I_F I_{\Theta^{-1}(G) }|| \mathcal{G} \right) \right]\\
								&= E\left[ E\left(I_{X^{-1}(A)} || \mathcal{G} \right) I_F I_{\Theta^{-1}(G)} \right]\\
								&= E\left[ E\left(I_{X^{-1}(A)} I_F I_{\Theta^{-1}(G)}|| \Theta, \mathcal{G} \right)  \right]\\
								&= E\left[I_{X^{-1}(A)} I_F I_{\Theta^{-1}(G)} \right]\\
								&= P(X^{-1}(A) \cap \Theta^{-1}(G) \cap F)
					\end{align*}	


\end{proof}

\section{Equivalence of Frequentist \& Bayesian Sufficiency}\index{Equivalence of Frequentist \& Bayesian Sufficiency}


Here we have, \\

$(\Omega_\Theta, \mathcal{F}_\Theta, \mu_\Theta), (\Omega_X, \mathcal{F}_X, \mu_X),(\Omega_T, \mathcal{F}_T)$\\

Where 

		$$T: \Omega_X \rightarrow \Omega_T \textcircled{m} \mathcal{F}_X / \mathcal{F}_T $$

is called a statistic. 

		$$T = T(X) \text{ or } T \circ X = T(X(\omega)) $$

In fewquentist setting, we say that T is \textbf{suffienct} if $P_{X|T, \Theta}$ does not depend on $\Theta$. It can be easily verified (see Homework) that $P_{X|T, \Theta}$ doesn't depend on $\Theta$ implies that

		$$P_{X|T, \Theta} = P_{X|T} \text{ a.s. P} $$

This is "nearly" fequentist. Above is exchangeable with "$X \indep \Theta| T$", but can't say this in frequentist setting. 


		$$P_{\Theta | T, X}  =  P_{\Theta|T} \Leftrightarrow P_{\Theta|X} = P_{\Theta|T}$$

That is to say that a statistic, T, is sufficient for $\Theta$ if and only iff the posterir distribution of $\Theta|X$ is the same as the posterior distribution of $\Theta|T$. This would be used in a Bayesian setting. 

\begin{definition}[Bayesian Sufficient]
	We say that $T \circ X$ is \textbf{Bayesian sufficient} if 

			$$P_{\Theta | X}  =  P_{\Theta|T} \text{ a.s. P } $$
\end{definition}


\textbf{Lemma 1.1} (HW 2) Suppose that $f(\theta)$ is a p.d.f such that

		$$f(\theta) \propto exp\{-a\theta^2 + b\theta \}, \quad a > 0 $$

		Then, 

		\begin{enumerate}
			\item $\theta \sim N( \frac{b}{2a}, \frac{1}{2a})$
			\item $\int exp\{-a\theta^2 + b\theta \} d\theta =\sqrt{\frac{\pi}{a}} exp\{\frac{b^2}{4a} \} $
		\end{enumerate}


\begin{example}
	Suppose that

			$$X | \Theta \sim N(\Theta, \sigma^2) $$
			$$\Theta \sim N(\mu, \tau^2)$$

	Find $\pi_{\Theta|X} (\theta|x), f_X(x)$. 

	\textbf{Solution:}\\

	\begin{align*}
		\pi(\theta|x) &\propto f(x|\theta) \pi(\theta)\\
				&= \frac{1}{\sqrt{2 \pi \sigma^2}} exp\{-\frac{1}{2} \frac{(x-\theta)^2}{\sigma^2}\} * \frac{1}{\sqrt{2 \pi \tau^2}} exp\{-\frac{1}{2} \frac{(\theta -\mu)^2}{\tau^2}\}\\
				&\propto  exp\{-\frac{1}{2} \frac{(x-\theta)^2}{\sigma^2}\} *  exp\{-\frac{1}{2} \frac{(\theta -\mu)^2}{\tau^2}\}\\
				&= exp\{-(\frac{1}{2\sigma^2} + \frac{1}{2 \tau^2})\theta^2 + (\frac{x}{\sigma^2} + \frac{\theta}{\tau^2})\theta\}
	\end{align*}

Using  Lemma 1.1, 

		$$\theta|X \sim N\left(\frac{\frac{x}{\sigma^2} + \frac{\mu}{\tau^2}}{1/2(2\sigma^{-2} + 1/2\tau^{-1} )}, \frac{1}{\frac{1}{\sigma^2} + \frac{1}{\tau^2}}\right) $$

How about $f_X(x)$?\\

		\begin{align*}
			f_X(x) &= \int f(x|\theta) \pi(\theta) d\theta\\
					&\vdots\\
					&= \frac{1}{{2 \pi \sigma \tau}} * exp\{-\frac{x^2}{2\sigma^2} -  \frac{\mu^2}{2\tau^2}\} \int exp\{- (\frac{1}{2\sigma^2} + \frac{1}{2 \tau^2})\theta^2 + (\frac{x}{\sigma^2} + \frac{\mu}{\tau^2})\theta \}\\
					&= \dots * \sqrt{\frac{\pi}{\frac{1}{2\sigma^2} + \frac{1}{2 \tau^2}}} exp\{\frac{(\frac{x}{\sigma^2} + \frac{\mu}{\tau^2})^2}{4 (\frac{1}{2\sigma^2} + \frac{1}{2\tau^2})}\}
		\end{align*}

		We want to identify this as a p.d.f of x, so we can treat anything that is not x as a constant. Using elementary algebra we get...

		\begin{align*}
			&\propto exp \{- (\frac{x^2}{2(\tau^2 + \sigma^2)} + \frac{x \mu}{(\sigma^2 + \tau^2)}) \}
		\end{align*}

		Applying Lemma 1.1 for x and simplifying, 

				$$X \sim N(\mu, \tau^2 + \sigma^2) $$


\end{example}

This can be extended to multivariate setting, 2-sample setting, ANOVA setting, regression setting, etc. It is essential to all aspects of linear models. 




	




% %------------------------------------------------




% %----------------------------------------------------------------------------------------
% %	CHAPTER 4
% %----------------------------------------------------------------------------------------

% \chapterimage{chapter_head_1.pdf} % Chapter heading image

% \chapter{Random Variable}


% %----------------------------------------------------------------------------------------
% %	CHAPTER 5
% %----------------------------------------------------------------------------------------

% \chapterimage{chapter_head_1.pdf} % Chapter heading image

% \chapter{Convergence in Probability/Limit Theorem}


% %----------------------------------------------------------------------------------------
% %	CHAPTER 6
% %----------------------------------------------------------------------------------------

% \chapterimage{chapter_head_1.pdf} % Chapter heading image

% \chapter{Radon-Nikodym Derivative Theorem}

% %----------------------------------------------------------------------------------------
% %	CHAPTER 7
% %----------------------------------------------------------------------------------------

% \chapterimage{chapter_head_1.pdf} % Chapter heading image

% \chapter{Special Topics}


%----------------------------------------------------------------------------------------
%	BIBLIOGRAPHY
%----------------------------------------------------------------------------------------

% \chapter*{Bibliography}
% \addcontentsline{toc}{chapter}{\textcolor{ocre}{Bibliography}}
% \section*{Books}
% \addcontentsline{toc}{section}{Books}
% \printbibliography[heading=bibempty,type=book]
% \section*{Articles}
% \addcontentsline{toc}{section}{Articles}
% \printbibliography[heading=bibempty,type=article]

%----------------------------------------------------------------------------------------
%	INDEX
%----------------------------------------------------------------------------------------

\cleardoublepage
\phantomsection
\setlength{\columnsep}{0.75cm}
\addcontentsline{toc}{chapter}{\textcolor{ocre}{Index}}
% \printindex

%----------------------------------------------------------------------------------------

\end{document}